% %%%%%%%%%%%%%%%%%%%%%%%%%%%%%%%%%%%%%%%%%%%%%%%%%%%%%%%%%%%%%%%%%%%%%%%%%%%%%%%%%%%%%%%
% Conclusion
% %%%%%%%%%%%%%%%%%%%%%%%%%%%%%%%%%%%%%%%%%%%%%%%%%%%%%%%%%%%%%%%%%%%%%%%%%%%%%%%%%%%%%%%

\FancyHeadRight{\sectitle}

\renewcommand{\sectitle}{Conclusion}

\phantomsection
\chapter*{\sectitle}
\addstarredchapter{\sectitle}

\minitoc

\newcounter{concluCnt}
\setcounter{concluCnt}{0}
\renewcommand{\thesection}{\arabic{concluCnt}}

% %%%%%%%%%%%%%%%%%%%%%%%%%%%%%%%%%%%%%%%%%%%%%%%%%%%%%%%%%%%%%%%%%%%%%%%%%%%%%%%%%%%%%%%
% Conclu 1

\addtocounter{concluCnt}{1}
\section{TODO}

\lettrine{I}{n} this thesis, we addressed the central question of \colorbf{Mulberry}{how hardware and architectural mechanisms can be designed to efficiently support computation with sparse data, while ensuring scalability, adaptability, and high performance across diverse workloads}.
Our approach focused on the challenges posed by irregular memory access patterns in sparse matrix computations, particularly in the context of the \textbf{\acrfull{op} \acrfull{spmv}} kernel.
We identified that the \acrshort{op} algorithm, while promising for parallelism and optimization with \acrshort{mv} kernels, generates ``random'' \glspl{reduction} (irregular updates on the output vector) that stress traditional memory systems.
To mitigate these challenges, we proposed \textbf{\acrlong{spdc}}, a novel \textbf{region-based \gls{redC}} architecture designed to optimize memory transactions for sparse data by leveraging the structural properties of \gls{banded} matrices.

Through this work, we developed a \textbf{lightweight, flexible solution} that can be integrated into any kernel suffering from ``random'' \glspl{reduction}, recognizing that caches are a fundamental component of memory hierarchies in CPUs, GPUs, and accelerators.
Our design prioritizes \textbf{scalability}, ensuring that the solution remains effective across different system scales.
Additionally, we addressed the critical influence of \textbf{matrix structure} on performance, conducting a detailed analysis of cache behavior specifically for \gls{banded} matrices.
By adapting our architecture to better handle dense and sparse regions, we provided insights into how matrix structure can inform the design of efficient architectures, with \textbf{\acrlong{spdc}} serving as a practical example.
This work also demonstrated how \textbf{simple yet targeted modifications} to cache design can yield significant performance improvements for \gls{reduction}-based kernels.

While our approach has limitations, such as the absence of a full evaluation on multi-level, multi-core systems and the early-stage refinement of \acrlong{spdc}, we intentionally focused on \textbf{precision and conceptual clarity} over raw performance metrics.
This emphasis on intrinsic understanding, rather than merely achieving impressive results, offers a more meaningful response to our central question and addresses a gap often overlooked in existing accelerator research.

\subsection*{Summary of Contributions}

\paragraph*{Chapter~\ref{cha:background}. Analysis of Sparse Data Computing}

We established a foundational understanding of sparse matrices, their storage formats, and the hardware challenges associated with their computation.
Our analysis of \textbf{\acrfull{ip}, \acrfull{op}, and \acrfull{gust}} algorithms revealed that \acrshort{op} and \acrshort{gust} are the most promising for hardware acceleration, particularly due to their potential for parallelism and data reuse.
We identified \textbf{\gls{gather} (\gls{intersection} searches)} and \textbf{\gls{scatter} (``random'' \glspl{reduction})} as the two core challenges in sparse matrix processing, which guided our subsequent architectural choices.

\paragraph*{Chapter~\ref{cha:soa}. Survey and Classification of State-of-the-Art Accelerators}
We conducted a comprehensive review of recent hardware accelerators for sparse matrix computations, categorizing them based on **autonomy, memory hierarchy position, algorithm compatibility, and hardware support for gather/scatter operations**. Our comparative analysis highlighted that **Gustavson-based accelerators** often achieve the best performance due to their balanced approach to input reuse and output generation. This survey informed the design of **SpDCache**, positioning it as a **processor-assist, near-cache accelerator** optimized for **OP SpMV kernels**.

\paragraph*{Chapter~\ref{cha:spdc}. Design of SpDCache for "Random" Reductions}
We introduced **SpDCache**, a **region-based reduction cache** that partitions memory into **In-Band (IBR) and Out-of-Band (OBR) regions** to optimize handling of both dense and sparse data. Our **Python-based simulation model** demonstrated that SpDCache significantly reduces memory traffic compared to generic caches, particularly for **banded matrices**. The architecture’s **virtual partitioning** and **reduction support** enable efficient management of irregular updates, addressing the scatter problem inherent in OP algorithms.

\paragraph*{Chapter~\ref{cha:theory}. Probabilistic Modeling and Exploration of Cache Behavior}
We developed a **hybrid probabilistic and analytical model** to simulate SpDCache’s behavior during OP SpMV operations. This model allowed us to explore how different cache configurations impact memory traffic, particularly for **banded matrices**. Our analysis revealed **two distinct groups of matrices** based on their **fine-grain structure**, as measured by the **$d_{CD}$ metric** (density of non-zero elements in the central diagonal region). This distinction underscores the importance of **adapting cache regions** to the specific structure of the input matrices.

\paragraph*{Chapter~\ref{cha:arch}. Microarchitecture and RTL Evaluation}
We detailed the **microarchitectural implementation of SpDCache**, featuring **three dedicated regions: Generic Region Cache (GRC), Dense Region Cache (DRC), and Sparse Region Table (SRT)**. Our **RTL simulations** validated the architecture’s effectiveness, showing **significant reductions in memory traffic** and **improved cache utilization**. The results confirmed that SpDCache efficiently handles both **coarse-grain** (region-based partitioning) and **fine-grain** (matrix structure-dependent) optimizations, making it adaptable to a wide range of sparse matrices.

\paragraph*{Chapter~\ref{cha:optim}. Optimizations and Future Directions}
We explored **RTL-level optimizations** to refine SpDCache’s performance, including **multi-core integration** and **extensions to SpMSpM kernels**. These enhancements position SpDCache as a **scalable, versatile solution** for modern sparse matrix processing systems. Our findings lay the groundwork for future research in **high-performance, energy-efficient architectures** for sparse data computations, with a focus on **bandwidth efficiency** and **adaptability to diverse workloads**.

---

\subsection*{Key Results and Insights}

- **Algorithmic Choice:** OP and Gust algorithms are the most effective for hardware acceleration, with **Gustavson offering the best trade-off** for SpMSpM kernels.
- **Cache Optimization:** SpDCache’s **region-based partitioning** and **reduction support** significantly reduce memory traffic, particularly for matrices with **structured sparsity** (e.g., banded matrices).
- **Fine-Grain Structure Matters:** The **$d_{CD}$ metric** revealed two distinct matrix groups, highlighting the need for **adaptive cache configurations** tailored to matrix structure.
- **Scalability and Usability:** SpDCache’s **multi-core extensions** and **optimizations for SpMSpM** demonstrate its potential for **scalable, high-performance sparse matrix processing**.

---

\subsection*{Future Work}

While SpDCache advances the state-of-the-art in cache architectures for sparse data, several avenues remain for future exploration:
- **Multi-Level Cache Hierarchies:** Extending SpDCache to **multi-level cache systems** could further improve performance for large-scale applications.
- **Dynamic Region Sizing:** Implementing **runtime-adaptive region sizing** based on matrix structure could enhance flexibility and efficiency.
- **Integration with Emerging Technologies:** Exploring SpDCache’s compatibility with **in-memory computing** and **processing-in-memory (PIM)** architectures may unlock new performance gains.
- **Broader Application Domains:** Evaluating SpDCache in **graph analytics, deep learning, and scientific computing** could demonstrate its versatility across diverse sparse workloads.

---

\subsection*{Final Thoughts}

This thesis contributes a **novel, reduction-aware cache architecture** that bridges the gap between theoretical insights and practical hardware implementations for sparse matrix computing. By addressing the challenges of **irregular memory accesses** and **efficient data reuse**, SpDCache provides a **scalable, adaptable solution** for modern computational demands. As sparse data continues to grow in importance across scientific, machine learning, and graph analytics applications, architectures like SpDCache will play a critical role in enabling **high-performance, energy-efficient processing** of these challenging workloads.


%%%%%%%%%%%%%%%









Our contributions include:
- **A comprehensive analysis of sparse matrix algorithms and hardware challenges**, highlighting the inefficiencies of traditional caches in handling irregular accesses and reductions.
- **The design and evaluation of SpDCache**, which partitions memory into **In-Band (IBR) and Out-of-Band (OBR) regions** to efficiently manage both dense and sparse data, significantly reducing memory traffic compared to generic caches.
- **Theoretical and probabilistic modeling** to analyze cache behavior, revealing the importance of **fine-grain matrix structure** (measured by the $d_{CD}$ metric) in optimizing performance.
- **Microarchitectural implementation and RTL simulations**, validating SpDCache’s effectiveness in reducing memory traffic and improving cache utilization, while identifying two distinct groups of matrices with different performance profiles.
- **Optimizations and extensions** for multi-core systems and SpMSpM kernels, positioning SpDCache as a scalable solution for modern sparse matrix processing.

However, our work has some limitations:
- **Single-core, single-cache level focus**: Our evaluation was constrained to a single-core, single-cache level system, which limits raw performance compared to more complex, multi-level, or multi-core architectures.
- **Uniformity assumptions**: The analytical model assumes a uniform distribution of non-zero elements, which may not always reflect real-world matrix structures, potentially affecting performance predictions.
- **Hardware complexity**: The fully-associative nature of the Sparse Region Table (SRT) and specialized eviction policies introduce hardware complexity, which could impact scalability and implementation costs.
- **RTL implementation optimization**: The current RTL implementation is a first iteration and could be further refined, as some features are still missing or not fully optimized.

While SpDCache advances the state-of-the-art in cache architectures for sparse data, future work could explore **multi-level cache hierarchies, dynamic region sizing, and integration with emerging technologies** like in-memory computing to further enhance performance and scalability. Additionally, broader application domains, such as graph analytics and deep learning, could be investigated to demonstrate SpDCache’s versatility across diverse sparse workloads.


% \subsection{A Section}

% \subsubsection{A Subsection}

% \lipsum[3]

% %%%%%%%%%%%%%%%%%%%%%%%%%%%%%%%%%%%%%%%%%%%%%%%%%%%%%%%%%%%%%%%%%%%%%%%%%%%%%%%%%%%%%%%
% Conclu 2

% \addtocounter{concluCnt}{1}
% \section{Second element}

% \lipsum[1-5]

% %%%%%%%%%%%%%%%%%%%%%%%%%%%%%%%%%%%%%%%%%%%%%%%%%%%%%%%%%%%%%%%%%%%%%%%%%%%%%%%%%%%%%%%

\renewcommand{\thesection}{\arabic{chapter}.\arabic{section}}

\cleardoublepage %% Comment this line if this section occupies the last pages
% \adjustmtc %% Commented for the appendices minitoc to not disappear